\documentclass[10pt,a4paper,twoside,fancy,bg=print,twocolumn,openany,nodeprecatedcode]{dndbook}
\usepackage[utf8]{inputenc}
\usepackage{tabulary}
\usepackage[normalem]{ulem}
\usepackage{float}
\usepackage{csquotes}
\raggedbottom
\extrafloats{2000}
\setlist[description]{nolistsep,listparindent=3pt,topsep=6pt}
\begin{document}
\chapter*{Overminds}

\begin{DndReadAloud}
\enquote{Careful, Xorannox}, the queen's consort cautioned. \enquote{You tread heavily.}

\enquote{Tread heavily? My dear consort}, the Tyract's grin twisted, \enquote{I never touch the ground.}
\end{DndReadAloud}

Formally known by their endonym, vaurath (plural vaurathi), overminds plot and scheme against their chief rivals, the synlirii and the olothec, for control of the World Below.

\textbf{\textit{Psionic Eyes.}} An overmind is an enormous floating brain with a large central eye surrounded by many smaller embedded eyes. Even more alien, several psionic eyes float within inches of their head, each connected to a small brain that can manifest a unique psionic effect.

\textbf{\textit{Intelligent Loremasters.}} Like their rivals, overminds are loremasters of supreme intellect. They aren't usually wizards, but they all view the collection of knowledge and arcane sciences as the best tool for ultimate conquest. Overminds are cruel and capricious but enjoy tests of intellectual might. Their one vanity is their unshakable belief in their own intellectual superiority over all others.

\textbf{\textit{Overmind Lairs.}} Thanks to their innate flight, overminds prefer to build (or rather, have their thralls build) towers with no doors or entrances anywhere near the bottom.

\textbf{\textit{Rivalries and Negotiations.}} Unlike their rivals, overminds have an intense hatred of their own kind and never work together. However, overminds often enjoy diplomacy with other species, seeking to form alliances and build secret networks of agents throughout the World Below.

\textbf{\textit{Smelly Eggs.}} Overmind procreation involves one aberration laying a single egg in a slime pool and leaving it behind. When another overmind later detects the distinct smell of the egg, they spray their inseminating fluid into the pool. These actions are compelled by biological necessity, a compulsion even the overminds can't ignore.

Unwitting explorers sometimes accidentally abscond with an egg. If unfertilized, its bearer is likely to attract the attention of other overminds and synlirii that use the eggs in genetic experiments to create psionic creatures. If fertilized, the explorer could find themselves as a newly hatched overmind's first victim.

\subsection{Xorannox the Tyract}

\textbf{Xorannox} (\textit{ZOR-ah-nocks}) rules as Lord of the White Tower, a multi-level finger of alabaster stone. Commonly known as the Tyract—an ancient Deep Speech word that literally translates as "a king who rules with his teeth"—the overmind indeed consumes those who displease him.

Xorannox is chief of the Grasp, a secret organization made up of deep gnomes, dark elves, and humans. They seek to overthrow the voiceless talkers' great empire and place Xorannox above all, first as king, then as god.

The Tyract is a master strategist, always one step ahead of his enemies. Adventurers from the Mundane World may be surprised to discover their deep allies are members of the Grasp. When they meet Xorannox, he's delighted! He loves treating with humanoids! They have the same enemies, after all.

Unfortunately, no matter how useful or dependable a party of heroes might be, Xorannox is nearly incapable of resisting the urge to betray his allies. He assumes, as do most vaurathi, that the natural end of all alliances is betrayal. It's just a matter of who stabs who first.

\textbf{\textit{Eyes Always Shielded.}} Xorannox's psionic might is so great that unlike other overminds, his floating eyes are never without their psionic shield. When he is truly in distress, Xorannox can sacrifice one of his eyes to bolster his resistances.

\setlist[description]{nolistsep,listparindent=3pt,topsep=6pt,font={\bfseries\sffamily\small}}


\begin{DndMonster}[float*=b,width=\textwidth + 8pt]{Xorannox}
\begin{multicols}{2}
\DndMonsterType{Large aberration (Solo), lawful evil}
\DndMonsterBasics[
armor-class = {18 (natural armor)},
hit-points = {\DndDice{24d10 + 96}},
speed = {0 ft., fly 20 ft. \textit{(hover)}}
]
\DndMonsterAbilityScores[
 str = 10,
 dex = 16,
 con = 18,
 int = 19,
 wis = 14,
 cha = 18
]
\DndMonsterDetails[
saving-throws = {Int +9, Wis +7, Cha +9},
skills = {Arcana +9, Deception +9, Insight +7, Intimidation +9, Perception +7},
condition-immunities = {charmed, flanked, frightened, prone},
senses = {darkvision 120 ft., passive perception 17},
languages = {Common, Deep Speech, Undercommon},
challenge = 14,
]
\DndMonsterAction{Painful Resistance (3/Day)}
Eight eyes float around Xorannox and create his Eye Psionics. If Xorannox fails a saving throw, he can destroy a random floating eye (chosen by rolling a d8) and succeed instead. If an eye is destroyed, a new eye pops of out of Xorannox's face to replace it 24 hours later.

\DndMonsterSection{Actions}
\DndMonsterAction{Bite}
\textsl{Melee Weapon Attack:} 5 to hit, reach 5 ft., one target. \textsl{Hit:} 14 (4d6) piercing damage.

\DndMonsterAction{Eye Psionics}
Xorannox creates three of the following psionic eye effects. Unless otherwise stated, each eye targets a creature who he can see within 120 feet of him. He can't use the same effect twice on a turn.


\begin{description}
\item[1. Charm Beam.]
The targeted creature must succeed on a DC 17 Wisdom saving throw or be charmed by Xorannox for 1 hour, or until he harms the creature or one of their allies.

\item[2. Compulsion Beam.]
The targeted creature must succeed on a DC 17 Intelligence saving throw or use their reaction, if available, to move up to their speed toward the closest ally they can see and make one weapon attack against them. Creatures who can't be charmed are unaffected.

\item[3. Memory Beam.]
The targeted creature must make a DC 17 Intelligence saving throw. On a failed save, if the creature has unexpended spell slots, they expend a spell slot of their highest remaining level, with no effect. If they don't have unexpended spell slots, they instead lose proficiency with a weapon they're holding or carrying for 1 minute (save ends at end of turn).

\item[4. Toxic Vapors.]
Xorannox chooses a point he can see within 120 feet of him. Each creature within 10 feet of the point must succeed on a DC 17 Constitution saving throw or be poisoned for 1 minute (save ends at end of turn).

\item[5. Telekinetic Field.]
Xorannox chooses a point he can see within 120 feet of him. Each enemy within 10 feet of the point must succeed on a DC 17 Strength saving throw or be moved up to 30 feet in any direction and take 7 (2d6) force damage at the end of the move. Additionally, Xorannox can target any object within 10 feet of the point that weighs 300 pounds or less and isn't being worn or carried, moving it up to 30 feet in any direction. Xorannox can also exert fine control on objects with this effect, such as manipulating a simple tool or opening a door or a container.

\item[6. Lightning Bolt.]
Xorannox shoots a 5-foot-wide, 60-foot-long line of lightning. Each creature in the line must succeed on a DC 17 Dexterity saving throw, taking 36 (8d8) lightning damage on a failed save, or half as much damage on a successful one.

\item[7. Explosion.]
Xorannox chooses a point he can see within 120 feet of him. Each creature within 10 feet of the point must make a DC 17 Dexterity saving throw, taking 36 (8d8) fire damage on a failed save, or half as much damage on a successful one.

\item[8. Necrosis Beam.]
The targeted creature must succeed on a DC 17 Constitution saving throw or immediately take 55 (10d10) necrotic damage and lose 7 (2d6) hit points at the start of each of their turns. If this effect reduces a target's hit points to 0, they die. The effect ends if Xorannox dies or can't see the target. This psionic effect can't be used again while this effect persists.

\end{description}

\DndMonsterSection{Bonus Actions}
\DndMonsterAction{The Great Eye (Recharge 6)}
Xorannox's central eye turns solid black and projects a 150-foot cone of energy. If Xorannox or any creature in that area is affected by a spell, the spell's effects immediately end for that creature.

\DndMonsterSection{Reactions}
\DndMonsterAction{Cower!}
When a creature hits Xorannox with an attack, he shoots a psionic fear beam at them. The creature must succeed on a DC 17 Wisdom saving throw or take 21 (6d6) psychic damage and be frightened of Xorannox for 1 minute (save ends at end of turn).

\DndMonsterSection{Legendary Actions}
Xorannox has three villain actions. He can take each action once during an encounter after an enemy's turn. He can take these actions in any order but can use only one per round.

\begin{DndMonsterLegendaryActions}
\DndMonsterLegendaryAction{Action 1: Disruption Beams}{Xorannox shoots a psionic disruption beam at three creatures he can see within 120 feet of him. Each target must make a DC 17 Wisdom saving throw. On a failed save, whenever the target makes more than one weapon attack on a turn or casts a spell that isn't a cantrip, they take 14 (4d6) psychic damage. This effect lasts for 1 minute (save ends at end of turn).
}
\DndMonsterLegendaryAction{Action 2: Disappearing Act}{Xorannox turns invisible until the end of his next turn and teleports up to 120 feet to an unoccupied space he can see.
}
\DndMonsterLegendaryAction{Action 3: Megabeam}{Xorannox unleashes the effects of all of his remaining psionic eyes at once in a 20-foot-wide, 120-foot-long line. Each creature in that area must make a save against two random eye effects from Xorannox's Eye
}
\DndMonsterLegendaryAction{Psionics (reroll duplicates for each target)}{If the effect usually affects an area, it instead only affects individual creatures.
}
\end{DndMonsterLegendaryActions}
\end{multicols}
\end{DndMonster}


\end{document}
